%%%%%%%%%%%%%%%%%%%%%%%%%%%%%%%%%%%%%%%%%%%%%%%%%%%%%
%			HLAVIČKA								%
%%%%%%%%%%%%%%%%%%%%%%%%%%%%%%%%%%%%%%%%%%%%%%%%%%%%%
\documentclass[openany,12pt]{memoir}

\usepackage[utf8]{inputenc} 
\usepackage[czech]{babel}
\usepackage[T1]{fontenc}
\usepackage[top=1.5cm, bottom=2cm, left=2cm, right=2cm]{geometry}  % --> NASTAVENÍ OKRAJŮ
\usepackage{fancyhdr}
\usepackage{graphicx}
%\usepackage{lmodern}
\usepackage{tgschola}  % --- Silnější font
\usepackage{xwatermark}
\usepackage{xcolor}
\usepackage{changepage}
\usepackage{newunicodechar}



%%%%%% Package na zpěvník
\usepackage[full]{leadsheets}%http://mirrors.nic.cz/tex-archive/macros/latex/contrib/leadsheets/leadsheets_en.pdf   --> dokumentace	
\definesongtitletemplate{empty}{} 
\setchords{
format = \bfseries,   %tučné akordy
minor = {mi},% 
input-notation = {german},%
output-notation = {german}%
}
\definesongtitletemplate{empty}{} 

\newlength{\drop}
\newwatermark[allpages,color=red!50,angle=0,scale=2, xpos=0,ypos=0]{\includegraphics[width=5cm]{obr/pozadi2.jpg}} %--> dvojka na pozadí


%%%% Vlastní příkazy
\newcounter{Slokočet}   %Automatické číslování slok
\newcommand{\mezera}{\vspace*{0.5cm}}   %Horizontální odsazení slok
\newcommand{\stred}{5.2cm}   %%% Na zarovnání slok doprostřed, pozn. automatičtější zarovnávání na střed nejde
\newcommand{\refren}{\mezera \noindent \textbf{R:} } %refrén
\newcommand{\sloka}{\mezera \noindent \addtocounter{Slokočet}{1} \arabic{Slokočet}. } 	%sloka, která se automaticky čísluje
\newcommand{\ssloka}{\mezera \noindent}  % vlastní číslo sloky
\newcommand{\carka}{,\:}
\newcommand{\predehra}{\mezera \hspace*{-1.85cm} \textbf{Předehra:} }
\newcommand{\m}[1]{\color{white}{#1}}  %Pro akordy
\newcommand{\ap}{'} %Pro apostrof
\newcommand{\elipsa}{\kern\fontdimen3\font} %Příkaz pro lepší zacházení s výpustkami (=...); je to vpodstatě jen mezera mezi tečkama výpustky

%%%%%%%%%%%%%%%%%%%%%%%%%%%%%%%%%%%%%%%%%%%%%%%%%%%%%%
%			NÁVOD									 %
%%%%%%%%%%%%%%%%%%%%%%%%%%%%%%%%%%%%%%%%%%%%%%%%%%%%%%
%1. Věci v hlavičce IGNOROVAT
%2. Píseň psát do prostoru mezi \begin{song} a \end{song}
%3. další řádek se značí dvěma odsazeníma (= dvakrát stisknout enter)
%4. \refren vždy na začátku refrenu a \sloka na začátek sloky (automaticky se čísluje)
%  \ = alt gr + q ; [/] = alt gr f/g ; {/} = alt gr + b/n; ^ = alt gr + 3 + mezera
%Cokoliv napíšete do ^{  } se bude brát jako akord
%když se toto bude dotýkat nějakého slova (nebude mezi tím a slovem mezera)
%tak se akordy zjeví nad slovem, ale pište to před slovo
%Když se to nedotýká slova, tak akord lítá ve vzduchu a vytiskne se větší mezera
%První možnost je asi preferovanější
%5. Akordy stačí psát jen do první sloky, když se nezmění -- kytaristi to zvládnou



%\usepackage{subfiles}
%</++>

\begin{document}

\pagestyle{empty}
%\addtocontents{toc}{\protect\thispagestyle{empty}}
%\tableofcontents \thispagestyle{empty}\newpage
\newgeometry{top=1.5cm, bottom = 0cm, left = 2cm, right = 2cm}

\newpage

%Pro kompilaci je důležitý, 	aby ten input a \end{document} zůstali na posledních
%dvou řádcích!!!!
\begin{song}{title=\predtitle\centering Letím hledat světy \\\large Kašpárek v rohlíku \vspace*{-0.3cm}}  %% sem se napíše jméno songu a autor
\begin{centerjustified}

\sloka
	Mám ^{Dsus2\,\,}doma ^{{\color{white}\_\_\_\_}B}mimozemšťana ^{{\color{white}\_\_\_\_}Dsus2}schovanýho pod ^{{\color{white}\_}B}postelí 

   	spad' k nám ^{G}z nebe do zahrady, řek' mi ^{F{\color{white}\_\,}},,Ahoj, tak jsem ^{A{\color{white}\_\_}}tady.``

   	Spolu ^{Dsus2{\color{white}\_\_\_}}prohledáme ^{B{\color{white}\_\_\,}}světy, tiše ^{Dsus2{\color{white}\_\_\_}}vklouznem' do ^{B}rakety 

   	a jen ^*{G}má váme, máváme ^{F{\color{white}\_}}všem, který cestou ^*{A}po tkáme, jo ^{Dmi\z}hó\elipsa.\elipsa.\elipsa. 


\refren
 	^{Dmi{\color{white}\_}}Letim, letim ^{B{\color{white}\_\_}}hledat světy, ^{G{\color{white}\_\_\_}}slunce, hvězdy ^{F{\color{white}\_\_\_}}planety 

   	^{Dmi{\color{white}\_}}Najdu místo ^{B}na obloze, ^{G}kde parkujou ^*{A}ko mety 

  	 ^{Dmi{\color{white}\_}}Letim, letim ^*{B}hl edat světy, ^{G\,\,}chci bejt zpátky ^{F}do pěti 

  	 ^*{G}Le tim hledat světy, ^{A{\color{white}\_\_}}můžeš se mnou jet i ty 

\sloka
	Tam na ^{Dsus2}konci ^{{\color{white}\_}B}galaxie, tam to ^{Dsus2}lítá, tam to ^{B{\color{white}\_}}žije 

  	 tam je ^*{G}ka ždej druhej týpek ^{{\color{white}\_}F}zelenej jak ^{{\color{white}\_\_\_}A}pistácie 

  	 Co nás ^{Dsus2}čeká za ^{B}legraci ve ^{{\color{white}\_}Dsus2}vesmírný ^{{\color{white}\_}B}restauraci 	

  	 ^*{G}či šníky tu nemají, ^*{F}ta líře samy ^*{A}lít ají 

  	 ^{Dsus2\z}Už~jsme v půlce ^{B{\color{white}\_\_\_\_\_}}nekonečna, ^{Dsus2}je fakt hustá ^*{B}dr áha Mléčná

  	 ^*{G}ba cha, uber do zatáčky, ^*{F}na razíme ^{A}do šlehačky. ^{Dmi}Jau\elipsa.\elipsa.\elipsa.


\refren


 \refren
 	^{Dmi{\color{white}\_}}Letim, letim ^{B{\color{white}\_\_}}hledat světy, ^{G\z}slunce, hvězdy ^{F{\color{white}\_\_\_}}planety 

   	^{Dmi{\color{white}\_}}Najdu místo ^{B}na obloze, ^{G}kde parkujou ^*{A}ko mety 

  	^{Dmi{\color{white}\_}}Letim, letim ^*{B}hl edat světy, ^{G\,\,}chci bejt zpátky ^{F}do pěti 

   	^*{G}Lí tat, to je psina, teď už ^{A{\color{white}\_\_}}chápu Gagarina 


\end{centerjustified}
\setcounter{Slokočet}{0}
\end{song}

\end{document}
